\chapter{Introduction}

\section{The Arami Uplands}
The Arami Uplands occupy about 4,800~km\textsuperscript{2} of dissected
plateau between the Ireke River in the south and the Asiri River in the
north. The land rises from about 350~m along the Ireke valley to 900~m on
the Aramiwo Escarpment, the highest ground in the region. The main
settlements are Okeri, the district headquarters and seat of the Okeri
Forestry Research Institute, and Idofin, a market town on the northern
plateau.

\section{Climate}
The Uplands have two wet seasons, from April to July and from September
to October, separated by a short dry spell in August. The long dry
season runs from November to March and is broken by the dry, dusty
harmattan wind from the north. Mean annual rainfall falls from about
1,600~mm in the Ireke valley to about 1,200~mm on the northern plateau.

\section{Vegetation}
The region lies in the transition between the Guineo-Congolian rain
forest and the Sudanian savanna woodlands described by
\citet{white1983}. Three broad vegetation types are recognised in this
Flora.

\textbf{Lowland semi-deciduous forest} covers the southern slopes below
about 500~m. Its canopy trees include \taxon{Khaya grandifoliola},
\taxon{Entandrophragma cylindricum}, \taxon{Milicia excelsa} and
\taxon{Terminalia superba}.

\textbf{Forest--savanna mosaic} occupies the middle slopes between
about 500 and 700~m, where gallery forest along streams alternates
with grassland that burns most years.

\textbf{Guinea savanna woodland} covers the northern plateau above
about 700~m. Here the fire-tolerant \taxon{Vitellaria paradoxa},
\taxon{Parkia biglobosa} and \taxon{Pterocarpus erinaceus} are
characteristic, and many of them are kept in farmland because of the
value of their fruits and seeds.

\section{History of botanical collecting}
The first substantial collections from the Uplands were made by
E.~M.~Harding between 1934 and 1938, during forest reserve surveys
around Okeri. J.~A.~Olaniyan collected widely on the northern plateau
between 1961 and 1967 for the Savanna Botanical Station. The most
recent collections were made by the editor between 2019 and 2025 for
this Flora. Specimens from all three periods are cited in the species
treatments.

\section{Scope and sources}
This volume treats the economic trees of the Uplands. Descriptions
draw on the regional floras of \citet{hutchinson1954},
\citet{aubreville1959}, \citet{keay1989} and \citet{akeassi2001}, and
on the field guides of \citet{arbonnier2004} and
\citet{hawthorne2006}. Information on uses follows
\citet{burkill1985}, \citet{louppe2008} and \citet{orwa2009}.
Nomenclature follows the \textit{Shenzhen Code} \citep{turland2018},
and author abbreviations follow \citet{brummitt1992}.